% ===== SLIDE 5 — CNN =====
\begin{frame}{Redes Neuronales Convolucionales}
\framesubtitle{La tecnología detrás del diagnóstico}

\begin{columns}[T]
\begin{column}{0.5\textwidth}
Las CNN imitan la \alert{corteza visual humana} para procesar imágenes:

\begin{itemize}
    \item Extraen características \alert{automáticamente}
    \item Aprenden jerarquías de patrones
    \item No requieren reglas programadas manualmente
\end{itemize}
\end{column}
\begin{column}{0.45\textwidth}
\begin{block}{\faCogs\; Arquitectura}
\begin{enumerate}
    \item Entrada (píxeles)
    \item Convolución + ReLU
    \item Pooling (reducción)
    \item Capas totalmente conectadas
    \item Salida (Softmax)
\end{enumerate}
\end{block}
\end{column}
\end{columns}

\end{frame}
